\section{Exemplarische Konfigurationsdateien des Prototyps}
\label{app:configuration_examples}

Abschnitt~\ref{sec:implementation_configuration} beschreibt die Trennung
zwischen ausführbarer Programmlogik und fach-, modell- beziehungsweise
zielspezifischer Konfiguration im Prototyp. Die nachfolgenden Ausschnitte
zeigen exemplarisch, wie diese Konfiguration im finalen Implementierungsstand
abgebildet ist.

Die vollständigen Konfigurationsdateien sind Bestandteil des mit der Arbeit
eingereichten Prototyps. Im Anhang werden nur diejenigen Ausschnitte
wiedergegeben, die für die in Kapitel~\ref{sec:implementation} beschriebene
Trennung zwischen Verarbeitungskern und Konfiguration relevant sind.
Die verwendeten Persona-Definitionen sind bereits vollständig in
Anhang~\ref{app:persona_definitions} dokumentiert.


\subsection{Zuordnung von Rollen und Personas}
\label{app:configuration_team}

Die Zusammenstellung der Persona-basierten Verarbeitung erfolgt über
strukturierte Teamkonfigurationen. Für die Interpretation quellengebundener
Engineering-Information wird im Prototyp die Datei
\texttt{source\_evidence\_interpretation\_team.json} verwendet.
Sie ordnet einer gemeinsamen Agentenrolle drei unterschiedliche
Persona-Definitionen zu.

Der folgende Ausschnitt zeigt die Zuordnung der Personas zu den jeweiligen
Markdown-Dateien und ihren Perspektiven.

\begin{lstlisting}[
basicstyle=\ttfamily\footnotesize,
breaklines=true,
columns=fullflexible,
keepspaces=true,
showstringspaces=false
]
{
  "team_id": "TEAM_SOURCE_EVIDENCE_INTERPRETATION",
  "team_name": "Source Evidence Interpretation Team",
  "task_name": "Interpret fixed source-grounded Evidence",
  "role_id": "ROLE_SOURCE_EVIDENCE_INTERPRETER",
  "role_file": "agents/roles/source_evidence_interpreter.md",
  "default_execution_mode": "multi_persona",
  "members": [
    {
      "member_id": "LEGACY_INTERPRETER_001",
      "agent_id": "AGENT_LEGACY_LITERAL_INTERPRETER",
      "persona_id": "PERSONA_LEGACY_LITERAL_INTERPRETER",
      "persona_file":
        "agents/personas/legacy_interpretation/literal_interpreter.md",
      "perspective": "literal_source_faithful"
    },
    {
      "member_id": "LEGACY_INTERPRETER_002",
      "agent_id": "AGENT_LEGACY_SYSTEMS_ENGINEERING_INTERPRETER",
      "persona_id":
        "PERSONA_LEGACY_SYSTEMS_ENGINEERING_INTERPRETER",
      "persona_file":
        "agents/personas/legacy_interpretation/systems_engineering_interpreter.md",
      "perspective": "systems_engineering"
    },
    {
      "member_id": "LEGACY_INTERPRETER_003",
      "agent_id": "AGENT_LEGACY_SKEPTICAL_AMBIGUITY_INTERPRETER",
      "persona_id":
        "PERSONA_LEGACY_SKEPTICAL_AMBIGUITY_INTERPRETER",
      "persona_file":
        "agents/personas/legacy_interpretation/skeptical_ambiguity_interpreter.md",
      "perspective": "skeptical_ambiguity"
    }
  ]
}
\end{lstlisting}

Die eigentliche Perspektive der Verarbeitung wird damit nicht innerhalb der
ausführbaren Programmlogik festgelegt. Die Teamkonfiguration verweist auf
separate Persona-Dateien, deren Inhalte in
Anhang~\ref{app:persona_definitions} wiedergegeben sind.


\subsection{Konfiguration des Modellrahmens}
\label{app:configuration_framework}

Der für die Modellbildung verwendete RFLP-basierte Modellrahmen liegt als
versionierte JSON-Konfiguration vor. Die Datei
\texttt{turing\_rflp\_framework.json} beschreibt die verfügbaren
Modellierungsebenen und die darin vorgesehenen Modellbereiche.

Der folgende Ausschnitt zeigt exemplarisch die Systemebene des
Modellrahmens.

\begin{lstlisting}[
basicstyle=\ttfamily\footnotesize,
breaklines=true,
columns=fullflexible,
keepspaces=true,
showstringspaces=false
]
{
  "node_id": "FW_LEVEL_SYSTEM",
  "mapping_key": "system_level",
  "name": "System Level",
  "node_type": "level",
  "parent_node_id": null,
  "mapping_target": false,
  "order": 2
},
{
  "node_id": "FW_SYSTEM_REQUIREMENTS",
  "mapping_key": "system.requirements",
  "name": "Requirements",
  "node_type": "framework_node",
  "parent_node_id": "FW_LEVEL_SYSTEM",
  "mapping_target": true,
  "order": 1
},
{
  "node_id": "FW_SYSTEM_FUNCTIONAL",
  "mapping_key": "system.functional",
  "name": "Functional",
  "node_type": "framework_node",
  "parent_node_id": "FW_LEVEL_SYSTEM",
  "mapping_target": true,
  "order": 2
},
{
  "node_id": "FW_SYSTEM_LOGICAL",
  "mapping_key": "system.logical",
  "name": "Logical",
  "node_type": "framework_node",
  "parent_node_id": "FW_LEVEL_SYSTEM",
  "mapping_target": true,
  "order": 3
},
{
  "node_id": "FW_SYSTEM_PHYSICAL",
  "mapping_key": "system.physical",
  "name": "Physical",
  "node_type": "framework_node",
  "parent_node_id": "FW_LEVEL_SYSTEM",
  "mapping_target": true,
  "order": 4
}
\end{lstlisting}

Die Struktur des Modellrahmens wird damit über stabile Identitäten und
explizite Beziehungen zwischen den Modellbereichen beschrieben. Das
ergänzende Modellstrukturprofil legt darauf aufbauend fest, welche
Elementtypen in den jeweiligen Modellbereichen zulässig sind.


\subsection{Konfiguration der SysML-v2-Erzeugung}
\label{app:configuration_generation}

Auch die Abbildung des internen Modellzustands auf die unterstützte
SysML-v2-Zielrepräsentation wird über eine versionierte Konfiguration
beschrieben. Die Datei
\texttt{turing\_sysml\_v2\_generation\_profile.json} enthält explizite
Abbildungsregeln für unterstützte und nicht unterstützte Modellinhalte.

Der folgende Ausschnitt zeigt eine Regel für die Abbildung einer Funktion
auf ein unterstütztes SysML-v2-Konstrukt.

\begin{lstlisting}[
basicstyle=\ttfamily\footnotesize,
breaklines=true,
columns=fullflexible,
keepspaces=true,
showstringspaces=false
]
{
  "rule_id": "J2_ELEMENT_006",
  "model_area": "system.functional",
  "element_type": "function",
  "mapping_status": "supported",
  "target_construct_id": "TN_006",
  "production_generation_allowed": true,
  "rationale":
    "An IEM function is an individual reviewed engineering behavior. It maps to an Action Usage/Feature so relationships can reference the generated occurrence without treating it as a reusable type.",
  "target_element_kind": "feature"
}
\end{lstlisting}

Nicht unterstützte Abbildungen werden im selben Profil explizit als solche
gekennzeichnet. Der Prototyp kann dadurch zwischen freigegebenen
Zielabbildungen und noch nicht unterstützten Modellinhalten unterscheiden,
ohne dafür eine implizite Ersatzabbildung innerhalb der Programmlogik
vorzunehmen.

Die exemplarischen Konfigurationen verdeutlichen damit die in
Abschnitt~\ref{sec:implementation_configuration} beschriebene Trennung:
Die grundlegenden Verarbeitungsmechanismen werden durch die Implementierung
bereitgestellt, während Perspektiven, Modellstruktur und Teile der
zielmodellspezifischen Abbildung über versionierte Konfigurationsdateien
festgelegt werden.